\chapter{द्विघात फलन और समीकरण}\label{ch:g11:quad}

रैखिक \omterm{def:g10:functions:function}{फलनों} के बाद सबसे सरल \omterm{def:g10:functions:function}{फलन} \omterm{def:g11:quad:quadratic}{द्विघात फलन} हैं, और सबसे पहले वही हैं जिनके
\omterm{def:g10:functions:graph}{आलेख} सचमुच वक्र होते हैं। यह अध्याय उनके लिए पूरा औज़ार-बक्सा तैयार करता है:
\omterm{thm:g11:quad:canonical}{मानक रूप}, \omterm{def:g11:quad:discriminant}{विविक्तकर}, किसी द्विघात व्यंजक का चिह्न, और \omterm{def:g11:quad:parabola}{परवलय} की ज्यामिति।

\section{मानक रूप}

\begin{definition}[द्विघात फलन]\label{def:g11:quad:quadratic}
\emph{द्विघात फलन}\index{द्विघात फलन} इस रूप का \omterm{def:g10:functions:function}{फलन} है
\[
f(x) = ax^2 + bx + c, \qquad a \neq 0,
\]
जहाँ $a$, $b$, $c$ \omterm{def:g10:numbers:sets}{वास्तविक संख्याएँ} हैं। व्यंजक $ax^2 + bx + c$ को \emph{द्विघात
त्रिपद} भी कहा जाता है।
\end{definition}

\begin{theorem}[मानक रूप]\label{thm:g11:quad:canonical}
हर \omterm{def:g11:quad:quadratic}{द्विघात फलन} को ठीक एक ही तरह से इस रूप में लिखा जा सकता है
\[
f(x) = a(x - \alpha)^2 + \beta,
\qquad\text{जहाँ } \alpha = -\frac{b}{2a} \text{ और } \beta = f(\alpha).
\]
यही $f$ का \emph{मानक रूप}\index{मानक रूप} है।
\end{theorem}

\begin{proof}
हम \emph{वर्ग पूरा करते हैं}। चूँकि $a \neq 0$ है, उसे पहले दो पदों में से बाहर
निकाल लीजिए:
\[
ax^2 + bx + c = a\left(x^2 + \frac{b}{a}x\right) + c .
\]
अब $x^2 + \frac{b}{a}x$ वर्ग $\left(x + \frac{b}{2a}\right)^2 = x^2 + \frac ba x + \frac{b^2}{4a^2}$ का आरंभ है, इसलिए $x^2 + \frac{b}{a}x = \left(x + \frac{b}{2a}\right)^2 -
\frac{b^2}{4a^2}$। प्रतिस्थापन करने पर:
\[
f(x) = a\left(x + \frac{b}{2a}\right)^2 - \frac{b^2}{4a} + c
= a\bigl(x - \alpha\bigr)^2 + \beta,
\]
जहाँ $\alpha = -\frac{b}{2a}$ और $\beta = c - \frac{b^2}{4a}$ हैं। $x = \alpha$ पर मान लेने से वर्ग वाला पद मर जाता है, इसलिए
$\beta = f(\alpha)$। अद्वितीयता: $a(x-\alpha)^2 + \beta$ का प्रसार करके $x$ के गुणांक मिलाने पर $\alpha = -\frac{b}{2a}$ आवश्यक
हो जाता है, और तब $\beta$ भी तय हो जाता है।
\end{proof}

\begin{example}\label{ex:g11:quad:canonical}
मान लीजिए $f(x) = 2x^2 - 8x + 3$ है। तब $\alpha = -\frac{-8}{2 \times 2} = 2$ और $\beta = f(2) = 8 - 16 + 3 = -5$, इसलिए
\[
f(x) = 2(x - 2)^2 - 5 .
\]
प्रसार करके जाँचिए: $2(x^2 - 4x + 4) - 5 = 2x^2 - 8x + 8 - 5 =
2x^2 - 8x + 3$।
\end{example}

\begin{proposition}[चरम मान और परिवर्तन]\label{prop:g11:quad:variations}
मान लीजिए $f(x) = a(x-\alpha)^2 + \beta$ है।
\begin{enumerate}
  \item यदि $a > 0$ हो, तो $f$ $\intoc{-\infty}{\alpha}$ पर \omterm{def:g10:functions:variations}{ह्रासमान} और $\intco{\alpha}{+\infty}$ पर \omterm{def:g10:functions:variations}{वर्धमान} है: $f$
        का \emph{\omterm{def:g10:functions:extrema}{न्यूनतम}} $\beta$ है, जो $x = \alpha$ पर प्राप्त होता है।
  \item यदि $a < 0$ हो, तो परिवर्तन उलट जाते हैं: $f$ का $x = \alpha$ पर
        \emph{\omterm{def:g10:functions:extrema}{अधिकतम}} $\beta$ है।
\end{enumerate}
\end{proposition}

\begin{proof}
मान लीजिए $a > 0$ है और $\alpha \leq u < v$ है। तब
\[
f(v) - f(u) = a\bigl((v-\alpha)^2 - (u-\alpha)^2\bigr)
= a\,(v - u)\,\bigl((v-\alpha) + (u-\alpha)\bigr).
\]
हर गुणनखंड धनात्मक है: $a > 0$, $v - u > 0$, और $(v-\alpha) + (u-\alpha) > 0$, क्योंकि $v > \alpha$ और $u \geq \alpha$ हैं।
इसलिए $f(v) > f(u)$: $f$ $\intco{\alpha}{+\infty}$ पर \omterm{def:g10:functions:variations}{वर्धमान} है। $u < v \leq \alpha$ के लिए अंतिम गुणनखंड ऋणात्मक है,
इसलिए $f(v) < f(u)$: $f$ वहाँ \omterm{def:g10:functions:variations}{ह्रासमान} है। अंत में सभी $x$ के लिए $f(x) - f(\alpha) = a(x-\alpha)^2 \geq 0$, इसलिए $\beta = f(\alpha)$
\omterm{def:g10:functions:extrema}{न्यूनतम} है। स्थिति $a < 0$ ऊपर वाली बात $-f$ पर लगाने से निकल आती है।
\end{proof}

\begin{example}
एक किसान के पास एक सीधी दीवार के सहारे आयताकार खेत घेरने के लिए $100$ मीटर
बाड़ है (दीवार के साथ बाड़ की आवश्यकता नहीं)। यदि $x$ चौड़ाई हो, तो घिरा
हुआ क्षेत्रफल $A(x) = x(100 - 2x) = -2x^2 + 100x$ है। यहाँ $\alpha =
-\frac{100}{2\times(-2)} = 25$ और $A(25) = 25 \times 50 = 1250$ हैं: क्षेत्रफल $25$ m चौड़ाई पर
सबसे बड़ा होता है, और उसका मान $1250$ m$^2$ है।
\end{example}

\section{मूल और गुणनखंडन}

\begin{definition}[मूल, विविक्तकर]\label{def:g11:quad:discriminant}
त्रिपद $ax^2+bx+c$ का \emph{मूल}\index{मूल} \omterm{def:g10:algebra:equation}{समीकरण} $ax^2 + bx + c = 0$ का हल है। त्रिपद का
\emph{विविक्तकर}\index{विविक्तकर} यह संख्या है:
\[
\Delta = b^2 - 4ac .
\]
\end{definition}

\begin{theorem}[$ax^2+bx+c=0$ को हल करना]\label{thm:g11:quad:roots}
\begin{enumerate}
  \item यदि $\Delta > 0$ हो, तो \omterm{def:g10:algebra:equation}{समीकरण} के दो भिन्न \omterm{def:g11:quad:discriminant}{मूल} हैं
        \[
        x_1 = \frac{-b - \sqrt{\Delta}}{2a},
        \qquad
        x_2 = \frac{-b + \sqrt{\Delta}}{2a},
        \]
        और $ax^2+bx+c = a(x - x_1)(x - x_2)$।
  \item यदि $\Delta = 0$ हो, तो \omterm{def:g10:algebra:equation}{समीकरण} का अकेला \omterm{def:g11:quad:discriminant}{मूल} $x_0 = -\frac{b}{2a}$ है, और $ax^2+bx+c = a(x - x_0)^2$।
  \item यदि $\Delta < 0$ हो, तो \omterm{def:g10:algebra:equation}{समीकरण} का कोई वास्तविक हल नहीं है, और त्रिपद का
        वास्तविक गुणांकों के साथ \omterm{def:g10:algebra:expand}{गुणनखंडन} नहीं हो सकता।
\end{enumerate}
\end{theorem}

\begin{proof}
\cref{thm:g11:quad:canonical} में निकाले गए \omterm{thm:g11:quad:canonical}{मानक रूप} से शुरू कीजिए:
\[
ax^2+bx+c
= a\left(x + \frac{b}{2a}\right)^2 + c - \frac{b^2}{4a}
= a\left[\left(x + \frac{b}{2a}\right)^2 - \frac{\Delta}{4a^2}\right],
\]
क्योंकि $c - \frac{b^2}{4a} = \frac{4ac - b^2}{4a} = -\frac{\Delta}{4a} =
a \times \left(-\frac{\Delta}{4a^2}\right)$ है।

\emph{1.} यदि $\Delta > 0$ हो, तो $\frac{\Delta}{4a^2} =
\left(\frac{\sqrt\Delta}{2a}\right)^2$ और कोष्ठक दो वर्गों का अंतर है:
\[
ax^2+bx+c = a\left(x + \frac{b}{2a} - \frac{\sqrt\Delta}{2a}\right)
\left(x + \frac{b}{2a} + \frac{\sqrt\Delta}{2a}\right)
= a(x - x_2)(x - x_1).
\]
गुणनफल ठीक तब शून्य है जब उसका कोई एक गुणनखंड शून्य हो, जिससे दोनों \omterm{def:g11:quad:discriminant}{मूल} मिल
जाते हैं, और वे भिन्न हैं क्योंकि $\sqrt\Delta \neq 0$ है।

\emph{2.} यदि $\Delta = 0$ हो, तो कोष्ठक $\left(x + \frac{b}{2a}\right)^2$ है, जो केवल $x_0 = -\frac{b}{2a}$ पर लुप्त होता है।

\emph{3.} यदि $\Delta < 0$ हो, तो $-\frac{\Delta}{4a^2} > 0$, इसलिए कोष्ठक एक वर्ग और एक धनात्मक संख्या
का योग है: वह कभी लुप्त नहीं होता, और \omterm{def:g10:algebra:equation}{समीकरण} का कोई वास्तविक हल नहीं है।
कोई \omterm{def:g10:algebra:expand}{गुणनखंडन} $a(x-x_1)(x-x_2)$ \omterm{def:g11:quad:discriminant}{मूल} पैदा कर देता, इसलिए ऐसा कोई \omterm{def:g10:algebra:expand}{गुणनखंडन} नहीं है।
\end{proof}

\begin{example}\label{ex:g11:quad:solving}
$2x^2 - 8x + 3 = 0$ हल कीजिए। यहाँ $\Delta = (-8)^2 - 4 \times 2 \times 3 =
64 - 24 = 40 > 0$, इसलिए दो \omterm{def:g11:quad:discriminant}{मूल} हैं:
\[
x_{1,2} = \frac{8 \pm \sqrt{40}}{4} = \frac{8 \pm 2\sqrt{10}}{4}
= 2 \pm \frac{\sqrt{10}}{2}.
\]
$x^2 + x + 1 = 0$ के लिए: $\Delta = 1 - 4 = -3 < 0$, कोई वास्तविक हल नहीं।
\end{example}

\begin{proposition}[मूलों का योग और गुणनफल]\label{prop:g11:quad:vieta}
यदि $\Delta \geq 0$ और $x_1, x_2$ \omterm{def:g11:quad:discriminant}{मूल} हों ($\Delta = 0$ होने पर बराबर), तो
\[
x_1 + x_2 = -\frac{b}{a},
\qquad
x_1 x_2 = \frac{c}{a} .
\]
इसके उलट, योग $s$ और गुणनफल $p$ वाली दो संख्याएँ $x^2 - sx + p = 0$ के हल होती हैं।
\end{proposition}

\begin{proof}
\cref{thm:g11:quad:roots} के \omterm{def:g10:algebra:expand}{गुणनखंडन} का प्रसार कीजिए:
\[
a(x - x_1)(x - x_2) = a x^2 - a(x_1 + x_2)x + a\,x_1 x_2 .
\]
गुणांकों को $ax^2 + bx + c$ से मिलाने पर $-a(x_1+x_2) = b$ और $a\,x_1x_2 = c$ मिलते हैं। विलोम के लिए: यदि
$u + v = s$ और $uv = p$ हों, तो $(x-u)(x-v) = x^2 - sx + p$, इसलिए $u$ और $v$ उसके \omterm{def:g11:quad:discriminant}{मूल} हैं।
\end{proof}

\begin{example}
योग $10$ और गुणनफल $21$ वाली दो संख्याएँ ज्ञात कीजिए। वे $x^2 - 10x + 21 = 0$ को हल
करती हैं; $\Delta = 100 - 84 = 16$, \omterm{def:g11:quad:discriminant}{मूल} $\frac{10 \pm 4}{2} = 7$ और $3$। \emph{\omterm{def:g11:quad:discriminant}{मूल} झट से पहचानना:} चूँकि $x^2 - 5x + 4$
के गुणांक $1 - 5 + 4 = 0$ को संतुष्ट करते हैं, इसलिए संख्या $1$ एक \omterm{def:g11:quad:discriminant}{मूल} है, और दूसरा
$\frac{c}{a} = 4$ है (उनका गुणनफल $\frac ca$ होना ही चाहिए)।
\end{example}

\section{द्विघात व्यंजक का चिह्न}

\begin{theorem}[त्रिपद का चिह्न]\label{thm:g11:quad:sign}
मान लीजिए $f(x) = ax^2+bx+c$ है, जहाँ $a \neq 0$ है।
\begin{enumerate}
  \item यदि $\Delta > 0$ हो और \omterm{def:g11:quad:discriminant}{मूल} $x_1 < x_2$ हों: $f(x)$ का चिह्न $\intcc{x_1}{x_2}$ के बाहर $a$
        जैसा है, और $\intoo{x_1}{x_2}$ पर $-a$ जैसा।
  \item यदि $\Delta = 0$ हो: $f(x)$ का चिह्न सभी $x \neq x_0$ के लिए $a$ जैसा है, और वह
        $x_0$ पर लुप्त होता है।
  \item यदि $\Delta < 0$ हो: $f(x)$ का चिह्न सभी $x$ के लिए $a$ जैसा है।
\end{enumerate}
संक्षेप में: \emph{त्रिपद का चिह्न $a$ जैसा होता है, सिवाय उसके \omterm{def:g11:quad:discriminant}{मूलों} के
बीच}।
\end{theorem}

\begin{proof}
\emph{1.} $f(x) = a(x-x_1)(x-x_2)$ लिखिए और गुणनखंडों के चिह्न देखिए। यदि $x < x_1$ हो: $x - x_1$ और
$x - x_2$ दोनों ऋणात्मक हैं, उनका गुणनफल धनात्मक है, इसलिए $f(x)$ का चिह्न $a$
जैसा है। यदि $x_1 < x < x_2$ हो: गुणनखंडों के चिह्न विपरीत हैं, गुणनफल ऋणात्मक है, और
$f(x)$ का चिह्न $-a$ जैसा है। यदि $x > x_2$ हो: दोनों गुणनखंड धनात्मक हैं और
$f(x)$ का चिह्न फिर $a$ जैसा है।

\emph{2.} $f(x) = a(x - x_0)^2$, और $x \neq x_0$ के लिए $(x-x_0)^2 > 0$।

\emph{3.} \cref{thm:g11:quad:roots} की उपपत्ति से $f(x) = a\left[\left(x + \frac{b}{2a}\right)^2 +
\left(-\frac{\Delta}{4a^2}\right)\right]$, और कोष्ठक सदा धनात्मक है।
\end{proof}

\begin{method}[द्विघात असमिकाएँ हल करना]\label{met:g11:quad:inequalities}
$ax^2 + bx + c \leq 0$ (या $\geq$, $<$, $>$) हल करने के लिए:
\begin{enumerate}
  \item $\Delta$ और, यदि हों तो, \omterm{def:g11:quad:discriminant}{मूल} निकालिए;
  \item \cref{thm:g11:quad:sign} का उपयोग करके \omterm{def:g10:algebra:signtable}{चिह्न-सारणी} बनाइए (\omterm{def:g11:quad:discriminant}{मूलों} के बाहर $a$ का
        चिह्न);
  \item हल-अंतराल पढ़ लीजिए, और ध्यान रखिए कि \omterm{def:g11:quad:discriminant}{मूल} स्वयं सम्मिलित हैं (चौड़ी
        असमिका) या नहीं (कड़ी)।
\end{enumerate}
\end{method}

\begin{example}\label{ex:g11:quad:inequality}
$-x^2 + 3x + 4 \geq 0$ हल कीजिए। यहाँ $\Delta = 9 + 16 = 25$, और \omterm{def:g11:quad:discriminant}{मूल} $\frac{-3 \pm 5}{-2}$ हैं, अर्थात् $x_1 = -1$ और $x_2 = 4$।
चूँकि $a = -1 < 0$ है, त्रिपद \omterm{def:g11:quad:discriminant}{मूलों} के बाहर ऋणात्मक और उनके बीच धनात्मक है। चौड़ी
असमिका के साथ हल-समुच्चय $\intcc{-1}{4}$ है।
\end{example}

\section{परवलय}

\begin{definition}[परवलय]\label{def:g11:quad:parabola}
\omterm{def:g11:quad:quadratic}{द्विघात फलन} $f(x) = a(x-\alpha)^2+\beta$ का \omterm{def:g10:functions:graph}{आलेख} शीर्ष $S(\alpha, \beta)$ वाला एक
\emph{परवलय}\index{परवलय} है। $a > 0$ होने पर वह ऊपर की ओर खुलता है, $a < 0$
होने पर नीचे की ओर, और वह ऊर्ध्वाधर रेखा $x = \alpha$ के परितः सममित है।
\end{definition}

\begin{remark}
सममिति की जाँच आसान है: किसी भी $h$ के लिए $f(\alpha + h) = ah^2 + \beta = f(\alpha - h)$, इसलिए रेखा $x = \alpha$ से
समान क्षैतिज दूरी पर पड़ने वाले दो बिंदुओं की ऊँचाई एक ही होती है।
\end{remark}

\begin{omfigure}
\begin{tikzpicture}
\begin{axis}[omaxis,
  width=5.1cm, height=4.8cm,
  xmin=-0.8, xmax=4.4, ymin=-1.7, ymax=3.6,
  xtick={1,3}, xticklabels={$x_1$,$x_2$}, ytick=\empty,
  title={$\Delta > 0$}]
  \addplot[omDef, thick, domain=-0.15:4.15] {(x-1)*(x-3)};
  \fill (axis cs:1,0) circle (1.5pt);
  \fill (axis cs:3,0) circle (1.5pt);
  \fill (axis cs:2,-1) circle (1.5pt) node[below, font=\small] {$S$};
\end{axis}
\end{tikzpicture}%
\begin{tikzpicture}
\begin{axis}[omaxis,
  width=5.1cm, height=4.8cm,
  xmin=-0.8, xmax=4.4, ymin=-1.7, ymax=3.6,
  xtick={2}, xticklabels={$x_0$}, ytick=\empty,
  title={$\Delta = 0$}]
  \addplot[omDef, thick, domain=0.2:3.8] {(x-2)^2};
  \fill (axis cs:2,0) circle (1.5pt)
    node[above right=1pt, font=\small] {$S$};
\end{axis}
\end{tikzpicture}%
\begin{tikzpicture}
\begin{axis}[omaxis,
  width=5.1cm, height=4.8cm,
  xmin=-0.8, xmax=4.4, ymin=-1.7, ymax=3.6,
  xtick=\empty, ytick=\empty,
  title={$\Delta < 0$}]
  \addplot[omDef, thick, domain=0.35:3.65] {(x-2)^2 + 0.8};
  \fill (axis cs:2,0.8) circle (1.5pt) node[below, font=\small] {$S$};
\end{axis}
\end{tikzpicture}

{\small $x$-अक्ष के सापेक्ष ऊपर की ओर खुलने वाले \omterm{def:g11:quad:parabola}{परवलय} ($a > 0$) की तीन
स्थितियाँ: दो \omterm{def:g11:quad:discriminant}{मूल}, एक द्विगुण \omterm{def:g11:quad:discriminant}{मूल}, कोई \omterm{def:g11:quad:discriminant}{मूल} नहीं। शीर्ष $S$ का \omterm{def:g10:coordgeom:system}{भुज} $\alpha = -\frac{b}{2a}$
है।}
\end{omfigure}

\begin{omfigure}
\begin{tikzpicture}
\begin{axis}[omaxis,
  width=8.6cm, height=5.6cm,
  xmin=-2.2, xmax=5.4, ymin=-4.4, ymax=6.2,
  xtick={-1,4}, xticklabels={$x_1$,$x_2$}, ytick=\empty]
  \addplot[name path=par, omDef, thick, domain=-1.75:4.75, samples=80]
    {-0.8*(x+1)*(x-4)};
  \path[name path=xaxis] (-1.75,0) -- (4.75,0);
  \addplot[omDef!20] fill between[of=par and xaxis,
    soft clip={domain=-1:4}];
  \fill (axis cs:-1,0) circle (1.5pt);
  \fill (axis cs:4,0) circle (1.5pt);
  \node[font=\small, omDef] at (axis cs:1.5,2.0) {$f > 0$};
  \node[font=\small, omThm] at (axis cs:-1.75,-1.3) {$f < 0$};
  \node[font=\small, omThm] at (axis cs:4.75,-1.3) {$f < 0$};
\end{axis}
\end{tikzpicture}

{\small दो \omterm{def:g11:quad:discriminant}{मूलों} वाले नीचे खुलने वाले त्रिपद ($a < 0$) का चिह्न: \omterm{def:g11:quad:discriminant}{मूलों} के बीच
$-a$ का चिह्न (छायांकित), बाहर $a$ का।}
\end{omfigure}

\begin{method}[परवलय पढ़ना]\label{met:g11:quad:reading}
$f(x) = ax^2+bx+c$ का रेखाचित्र बनाने के लिए: दिशा ($a$ का चिह्न), शीर्ष $\left(-\frac{b}{2a},\, f\!\left(-\frac{b}{2a}\right)\right)$,
$y$-अक्ष के साथ प्रतिच्छेद $(0, c)$, और यदि हों तो \omterm{def:g11:quad:discriminant}{मूल} ज्ञात कीजिए। फिर
सममिति-अक्ष $x = -\frac{b}{2a}$ हर निकाले हुए बिंदु को दो बार रख देता है।
\end{method}

\begin{example}
$f(x) = x^2 - 4x + 3$ के लिए: ऊपर की ओर, शीर्ष $(2, -1)$, $y$-अक्ष को $(0,3)$ पर काटता है, \omterm{def:g11:quad:discriminant}{मूल}
$1$ और $3$ (झट से पहचाने गए: $1 - 4 + 3 = 0$)। सममिति से बिंदु $(4, 3)$ भी वक्र पर
है, जो $(0, 3)$ का दर्पण-प्रतिबिंब है।
\end{example}

\section{अभ्यास}

\begin{exercise}[$\star$]\label{exo:g11:quad:1}
\omterm{def:g10:algebra:equation}{समीकरण} हल कीजिए:
\[
x^2 - 5x + 6 = 0, \qquad
2x^2 + 3x - 2 = 0, \qquad
x^2 + 2x + 5 = 0, \qquad
9x^2 - 6x + 1 = 0 .
\]
\end{exercise}

\begin{exercise}[$\star$]\label{exo:g11:quad:2}
\omterm{thm:g11:quad:canonical}{मानक रूप} में लिखिए और \omterm{def:g11:quad:parabola}{परवलय} का शीर्ष बताइए:
\[
f(x) = x^2 - 6x + 11, \qquad
g(x) = -2x^2 + 4x + 1, \qquad
h(x) = 3x^2 + 6x .
\]
\end{exercise}

\begin{exercise}[$\star$]\label{exo:g11:quad:3}
असमिकाएँ हल कीजिए:
\[
x^2 - 3x - 10 < 0, \qquad
2x^2 + x + 3 > 0, \qquad
-x^2 + 6x - 9 \geq 0 .
\]
\end{exercise}

\begin{exercise}[$\star$]\label{exo:g11:quad:4}
$\Delta$ निकाले बिना एक स्पष्ट \omterm{def:g11:quad:discriminant}{मूल} ज्ञात कीजिए, फिर दूसरा:
\[
x^2 - 7x + 6 = 0, \qquad
2x^2 + 5x - 7 = 0, \qquad
x^2 + 4x - 5 = 0 .
\]
(संकेत: $x = 1$ और $x = -1$ जाँचिए, फिर \omterm{def:g11:quad:discriminant}{मूलों} के गुणनफल का उपयोग कीजिए।)
\end{exercise}

\begin{exercise}[$\star\star$]\label{exo:g11:quad:5}
ऐसी दो संख्याएँ ज्ञात कीजिए जिनका योग $14$ और गुणनफल $45$ हो। फिर परिमाप
$34$ और क्षेत्रफल $60$ वाले आयत की विमाएँ ज्ञात कीजिए।
\end{exercise}

\begin{exercise}[$\star\star$]\label{exo:g11:quad:6}
मान लीजिए $f(x) = x^2 + (m-2)x + 1$ है, जहाँ $m$ एक वास्तविक प्राचल है। $m$ के किन मानों के
लिए \omterm{def:g10:algebra:equation}{समीकरण} $f(x) = 0$ का ठीक एक हल है? किनके लिए कोई हल नहीं?
\end{exercise}

\begin{exercise}[$\star\star$]\label{exo:g11:quad:7}
एक गेंद ऊपर की ओर फेंकी जाती है; $t$ सेकंड बाद उसकी ऊँचाई $h(t) = -5t^2 + 20t + 1$ (मीटर में)
है। \omterm{def:g10:functions:extrema}{अधिकतम} कितनी ऊँचाई तक वह पहुँचती है, और किस समय? किस समय-अंतराल में गेंद
कम से कम $16$ m ऊँचाई पर रहती है?
\end{exercise}

\begin{exercise}[$\star\star$]\label{exo:g11:quad:8}
$X = x^2$ रखकर \omterm{def:g10:algebra:equation}{समीकरण} $x^4 - 13x^2 + 36 = 0$ हल कीजिए (\emph{द्विवर्गीय} \omterm{def:g10:algebra:equation}{समीकरण})।
\end{exercise}

\begin{exercise}[$\star\star$]\label{exo:g11:quad:9}
$x$ के किन मानों के लिए \omterm{def:g11:quad:parabola}{परवलय} $y = x^2$ का बिंदु $M(x, x^2)$ रेखा $y = x + 2$ से कड़ाई से
नीचे है? रेखाचित्र पर व्याख्या कीजिए।
\end{exercise}

\begin{exercise}[$\star\star$]\label{exo:g11:quad:10}
किसी धनात्मक संख्या और उसके व्युत्क्रम का योग $\frac{13}{6}$ है। वह संख्या ज्ञात
कीजिए। (द्विघात \omterm{def:g10:algebra:equation}{समीकरण} पर लाइए।)
\end{exercise}

\begin{exercise}[$\star\star\star$]\label{exo:g11:quad:11}
मान लीजिए $\mathcal P$ \omterm{def:g11:quad:parabola}{परवलय} $y = x^2$ है और $d_m$ रेखा $y = mx - 1$ है, जहाँ $m$ एक
वास्तविक प्राचल है।
\begin{enumerate}
  \item $m$ के किन मानों के लिए $\mathcal P$ और $d_m$ दो बिंदुओं पर मिलते हैं?
        एक बिंदु पर? किसी बिंदु पर नहीं?
  \item जब वे ठीक एक बिंदु पर छूते हैं, तब स्पर्श बिंदु निकालिए। (\cref{ch:g11:deriv} में
        आप $d_m$ को उस बिंदु पर $\mathcal P$ की स्पर्श रेखा के रूप में पहचानेंगे।)
\end{enumerate}
\end{exercise}

\section{समस्या: परवलय का गुप्त बिंदु}

\begin{problem}[{सप्ताहांत समस्या --- नाभि और नियता: उपग्रह-तश्तरियाँ,
हेडलाइट और झूला पुल सब एक ही वक्र क्यों खींचते हैं, और वह वक्र मूलतः एक ही
है}]
\label{pb:g11:quad:1}
हर बालकनी की हर उपग्रह-तश्तरी अपने अभिग्राही को भीतर के एक निश्चित बिंदु पर
ताने रखती है; हर हेडलाइट अपना बल्ब उसी ख़ास जगह छिपाए रहती है। वह बिंदु ---
\emph{\omterm{pb:g11:quad:1}{नाभि}} --- वह रहस्य है जिसे \omterm{def:g11:quad:parabola}{परवलय} यूनानी ज्यामिति के समय से सँभाले हुए
है, और उसे बाहर निकालने के लिए दूरी सूत्र ही पर्याप्त है। रास्ते में यह
समस्या \omterm{def:g11:quad:discriminant}{विविक्तकर} का अभ्यास कराती है, दीवारों के ऊपर से प्रक्षेप्य फेंकती है,
एक पुल लटकाती है, और अंत में एक चकित कर देने वाला तथ्य सिद्ध करती है: ज़ूम
की छूट दे दें तो \emph{संसार में केवल एक ही \omterm{def:g11:quad:parabola}{परवलय} है}।

\medskip
\noindent\textbf{भाग I --- त्रिपद में प्रवाह।}
\begin{enumerate}
  \item $x^2 - 5x + 6 = 0$ हल कीजिए, फिर $2x^2 + x - 6 = 0$, फिर $x^2 + x + 1 = 0$ (\cref{thm:g11:quad:roots})।
  \item $2x^2 - 4x - 6$ को गुणनखंडित रूप में और \omterm{thm:g11:quad:canonical}{मानक रूप} में लिखिए (\cref{thm:g11:quad:canonical})। प्रसारित
        रूप के साथ अब आपके पास एक ही \omterm{def:g10:functions:function}{फलन} के तीन वेश हैं: बताइए कि हर वेश
        एक नज़र में कौन-सी विशेषता दिखा देता है (\omterm{def:g11:quad:discriminant}{मूल}, शीर्ष, $y$-अंतःखंड)।
  \item चिह्न के नियम (\cref{thm:g11:quad:sign}) से $-x^2 + 4x - 3 \geq 0$ हल कीजिए।
  \item वियेत जासूस की भूमिका में (\cref{prop:g11:quad:vieta}, जो \cref{pb:g10:algebra:1} में बेबीलोन को आगे
        बढ़ाता है): योग $7$ और गुणनफल $12$ वाली दो संख्याएँ उनका \omterm{def:g10:algebra:equation}{समीकरण}
        लिखकर ज्ञात कीजिए। फिर: $x^2 - 5x + c = 0$ का एक \omterm{def:g11:quad:discriminant}{मूल} $2$ है; कुछ भी हल किए
        बिना $c$ और दूसरा \omterm{def:g11:quad:discriminant}{मूल} ज्ञात कीजिए।
  \item $m$ के किन मानों के लिए $x^2 + mx + 9 = 0$ का द्विगुण \omterm{def:g11:quad:discriminant}{मूल} है? किनके लिए कोई
        \omterm{def:g11:quad:discriminant}{मूल} नहीं? किनके लिए दो \omterm{def:g11:quad:discriminant}{मूल}?
\end{enumerate}

\medskip
\noindent\textbf{भाग II --- \omterm{pb:g11:quad:1}{नाभि} का पर्दा उठता है।}
\omterm{def:g11:quad:parabola}{परवलय} $y = x^2$, बिंदु $F\left(0, \frac14\right)$ और \omterm{def:g10:algebra:equation}{समीकरण} $y = -\frac14$ वाली क्षैतिज रेखा $\delta$ पर काम
कीजिए।
\begin{enumerate}[resume]
  \item \omterm{def:g11:quad:parabola}{परवलय} के किसी बिंदु $P(x, x^2)$ के लिए $PF^2 = x^2 + \left(x^2 - \frac14\right)^2$ का प्रसार कीजिए और दिखाइए कि
        वह पूर्ण वर्ग $\left(x^2 + \frac14\right)^2$ है। इससे $PF$ निकालिए।
  \item $P$ की रेखा $\delta$ से दूरी निकालिए, और निष्कर्ष निकालिए:
        \emph{\omterm{def:g11:quad:parabola}{परवलय} का हर बिंदु $F$ से उतना ही दूर है जितना $\delta$ से}।
        ($F$ \emph{नाभि}\index{नाभि} है और $\delta$ \emph{नियता}\index{नियता}।)
  \item विलोम सिद्ध कीजिए: ऐसा बिंदु $P(x, y)$ जिसके लिए $PF = \operatorname{dist}(P, \delta)$ हो, $y = x^2$ को
        संतुष्ट करता ही है। (दोनों दूरियों का वर्ग कीजिए और सरल कीजिए।) इस
        तरह \omterm{def:g11:quad:parabola}{परवलय} एक \emph{बिंदुपथ} है, ठीक जैसे \cref{pb:g10:coordgeom:1} में वृत्त था: एक
        बिंदु से समान दूरी के बजाय एक बिंदु और एक रेखा से समान दूरी।
  \item चौड़े \omterm{def:g11:quad:parabola}{परवलय} $y = \frac{x^2}{4f}$ के लिए ($f > 0$ के साथ) प्रश्न 6 को ढालकर दिखाइए
        कि \omterm{pb:g11:quad:1}{नाभि} $(0, f)$ है और \omterm{pb:g11:quad:1}{नियता} $y = -f$। अनुप्रयोग: एक उपग्रह-तश्तरी $60$
        cm चौड़ी और $9$ cm गहरी है --- उसका किनारा $(30, 9)$ से होकर जाता है।
        $f$ ज्ञात कीजिए: अभिग्राही को शीर्ष से कितना ऊपर लटकाना चाहिए?
  \item तश्तरियाँ काम क्यों करती हैं: अक्ष के समांतर आने वाली सभी किरणें
        \emph{\omterm{pb:g11:quad:1}{नाभि} से होकर} परावर्तित होती हैं (और \omterm{pb:g11:quad:1}{नाभि} पर रखा बल्ब समांतर
        पुंज फेंकता है --- हेडलाइट उलटी चलती तश्तरियाँ ही हैं)। समझाइए कि
        प्रश्न 7 की समान-दूरी इसे प्रशंसनीय क्यों बना देती है (किरणों को
        आगे बढ़ते तरंगाग्र मानिए), और बताइए कि \cref{ch:g11:deriv} में आने वाला कौन-सा
        औज़ार (\cref{exo:g11:quad:11} देखिए) प्रशंसनीयता को उपपत्ति में बदल देता है।
\end{enumerate}

\medskip
\noindent\textbf{भाग III --- \omterm{def:g11:quad:parabola}{परवलय} काम पर।}
\begin{enumerate}[resume]
  \item एक गेंद पथ $y = x - \frac{x^2}{20}$ पर चलती है ($x$ और $y$ मीटर में)। ज्ञात कीजिए
        कि वह कहाँ गिरती है, उसके सबसे ऊँचे बिंदु का क्षैतिज स्थान क्या है,
        और \omterm{def:g10:functions:extrema}{अधिकतम} ऊँचाई कितनी है (\cref{prop:g11:quad:variations})।
  \item उसी पथ पर $x = 4$ पर एक दीवार खड़ी है। वहाँ गेंद की ऊँचाई निकालिए:
        क्या वह $3$ m की दीवार पार कर जाती है? $4$ m की?
  \item एक परवलयाकार मेहराब $40$ m चौड़ा है और उसकी \omterm{def:g10:functions:extrema}{अधिकतम} ऊँचाई $10$ m
        है: केंद्र पर आधारित \omterm{def:g10:coordgeom:system}{निर्देशांकों} में $y = 10 - \frac{x^2}{40}$। क्या $7$ m ऊँचा
        ट्रक केंद्र से $10$ m पर निकल सकता है? $15$ m पर?
  \item $p$ यूरो टिकट-दाम पर किसी रंगमंच की आय $R(p) = p\,(120 - 2p)$ है (दाम बढ़ने पर
        माँग घटती है)। आय को \omterm{def:g10:functions:extrema}{अधिकतम} करने वाला दाम और वह \omterm{def:g10:functions:extrema}{अधिकतम} आय ज्ञात
        कीजिए।
  \item गोल्डन गेट के मुख्य रस्से बहुत अच्छी यथार्थता के साथ परवलयाकार हैं
        (समान रूप से बँटा हुआ भार \omterm{def:g11:quad:parabola}{परवलय} बनाता है --- स्वतंत्र रूप से लटकती
        ज़ंजीर \emph{नहीं} बनाती: उसका वक्र, कैटेनरी, कक्षा 12 के चरघातांकी
        \omterm{def:g10:functions:function}{फलन} की माँग करता है)। $1\,280$ m की चौड़ाई और $143$ m के झोल के साथ
        रस्सा केंद्र से $y = 143\left(\frac{x}{640}\right)^2$ का अनुसरण करता है। $x = 320$ m पर रस्सा अपने सबसे
        नीचे वाले बिंदु से कितना ऊपर है?
\end{enumerate}

\medskip
\noindent\textbf{भाग IV --- रेखाएँ, स्पर्श रेखाएँ, और आख़िरी अचरज।}
\begin{enumerate}[resume]
  \item \cref{exo:g11:quad:11} पूरा कीजिए: बाहरी बिंदु $(0, -1)$ से होकर जाने वाली रेखाएँ $y = mx - 1$
        $y = x^2$ से एक \omterm{def:g11:quad:discriminant}{विविक्तकर} के चिह्न के अनुसार मिलती हैं। $(0, -1)$ से खींची
        जा सकने वाली दोनों स्पर्श रेखाएँ और उनके स्पर्श बिंदु ज्ञात कीजिए।
  \item अब \emph{भीतरी} बिंदु $(0, 1)$ से होकर जाने वाली रेखाएँ $y = mx + 1$:
        प्रतिच्छेद-समीकरण का \omterm{def:g11:quad:discriminant}{विविक्तकर} निकालिए और दिखाइए कि ऐसी हर रेखा
        \omterm{def:g11:quad:parabola}{परवलय} को दो बार काटती है। सीख: भीतर से कोई स्पर्श रेखा नहीं ---
        ठीक जैसे वृत्त के लिए।
  \item नाभीय जीवा: $y = x^2$ को \omterm{pb:g11:quad:1}{नाभि} से होकर जाने वाली क्षैतिज रेखा $y = \frac14$ से
        काटिए। जीवा की लंबाई कितनी है? सामान्य नियम जाँचिए (तश्तरियों के
        लिए याद रखने लायक़): $y = \frac{x^2}{4f}$ के लिए \omterm{pb:g11:quad:1}{नाभि} से होकर \omterm{pb:g11:quad:1}{नियता} के समांतर जाने
        वाली जीवा की लंबाई $4f$ होती है।
  \item आख़िरी अचरज: \omterm{def:g11:quad:parabola}{परवलय} $y = x^2$ पर ज़ूम $(x, y) \mapsto (kx, ky)$ लगाइए और दिखाइए कि \omterm{def:g10:functions:function}{प्रतिबिंब}
        ठीक $y = \frac{x^2}{k}$ है। निष्कर्ष निकालिए: हर \omterm{def:g11:quad:parabola}{परवलय} $y = \frac{x^2}{4f}$ इकाई \omterm{def:g11:quad:parabola}{परवलय} का ही
        बड़ा किया हुआ रूप है --- भिन्न-भिन्न अनुपातों वाले वृत्तों और
        दीर्घवृत्तों के विपरीत, \emph{सभी \omterm{def:g11:quad:parabola}{परवलय} समरूप हैं}।
  \item समापन: \omterm{def:g11:quad:parabola}{परवलय} का चित्र पाँच पंक्तियों में --- तीन वेशों वाला एक
        \omterm{def:g10:algebra:equation}{समीकरण} (प्रश्न 2), एक बिंदुपथ (प्रश्न 7--8), अभियंता का परावर्तक
        (प्रश्न 9, 10), प्रक्षेप्य का पथ और पुल का रस्सा (प्रश्न 11--15),
        और ज़ूम की छूट पर एक ही वक्र (प्रश्न 19)। हर एक पर एक वाक्य।
\end{enumerate}
\end{problem}
